\chapter{Acquisition de compétences d'ingénieur}
\label{chap:competences}

\textbf{Structure provisoire :} le référentiel officiel de compétences de
TELECOM Nancy n'a pas été retrouvé dans les espaces de travail. Les situations
ci-dessous sont factuelles, mais leurs intitulés et niveaux devront être alignés
sur ce référentiel avant la remise.

\section{Concevoir un système d'intelligence artificielle sous contraintes}

La réalisation du \gls{hmoe} m'a conduit à traiter conjointement apprentissage,
architecture logicielle et exploitation. La contrainte matérielle interdisait
de maintenir quatre modèles complets. J'ai donc partagé une base quantifiée et
spécialisé son comportement avec quatre adaptateurs \gls{qlora}. Ce choix ne
suffisait cependant pas à obtenir un agent fiable : il a fallu préparer des
conversations conservant les couples appel--résultat, aligner les fenêtres
d'entraînement et d'inférence, sérialiser le changement d'adaptateur et limiter
la surface d'outils présentée à chaque expert.

Cette situation m'a appris à évaluer un choix de modèle à l'échelle du système.
La consommation mémoire, la qualité des données, le routage, les contrats
d'outils et le comportement en cas de sortie malformée sont interdépendants.
Une bonne perte de validation ne remplace donc ni les tests d'intégration ni une
campagne agentique sur des scénarios représentatifs.

\section{Analyser et maîtriser un système complexe}

LANCE relie topologie réseau, outils protocolaires, orchestration multi-phase,
modèle de langage et infrastructure de test. Pour éviter qu'une erreur se
propage silencieusement, j'ai transformé plusieurs attentes implicites en
invariants vérifiables : prérequis de phase, couverture minimale de
reconnaissance, provenance d'une preuve, périmètre réseau autorisé et filtrage
des constats avant le rapport final.

\section{Expérimenter et valider}

Le harness d'évaluation m'a obligé à distinguer une affirmation plausible d'un
résultat démontré. J'ai travaillé sur une correspondance structurelle entre
constats et vérité terrain, sur la recomposition du niveau de preuve depuis les
appels d'outils, puis sur la couverture de chemins ordonnés. Les erreurs
d'outils, les tests négatifs concluants et l'absence de test reçoivent des
statuts différents. Cette distinction évite qu'une panne soit interprétée comme
une absence de vulnérabilité ou qu'un texte assuré soit crédité sans exécution.

J'ai également appris à séparer les indicateurs internes des conclusions
scientifiques. Les pertes d'entraînement servent à détecter une divergence ou à
choisir un checkpoint ; elles ne mesurent pas directement la capacité d'un
expert à terminer une phase. De même, l'instantané expérimental actuel constitue
surtout un test d'intégration du pipeline et du harness. Une comparaison
contrôlée du \gls{hmoe} avec plusieurs baselines reste nécessaire.

\section{Industrialiser et documenter}

Le passage du prototype au service a demandé une frontière explicite entre code
versionné, données, checkpoints et adaptateurs finaux. La synchronisation repose
sur une liste blanche, des empreintes et des limites de taille ; le serveur est
exposé par une \gls{api} compatible avec le reste de LANCE et supervisé comme un
service. Côté benchmark, la composition des scénarios, leur déploiement et leur
évaluation partagent des contrats versionnés. J'ai ainsi appris qu'une
fonctionnalité n'est réellement exploitable que si elle peut être déployée,
observée, testée et reproduite.

\section{Auto-évaluation}

Je considère avoir progressé dans la conception d'architectures hybrides, où un
modèle probabiliste reste encadré par des mécanismes déterministes. Je dois
encore approfondir deux points : construire une campagne d'ablation permettant
d'attribuer précisément les gains au routage des experts, et calibrer les
heuristiques du graphe sur des observations empiriques. Le niveau officiel de
chaque compétence restera à compléter à partir du référentiel de TELECOM Nancy.
